\documentclass[12pt,a4paper]{article}

\usepackage{amsmath,amssymb,amsthm}
\usepackage{mathtools}
\usepackage{ytableau}
\usepackage{tikz}
\usepackage{enumitem}
\usepackage[margin=2.5cm]{geometry}
\usepackage{hyperref}

% Theorem environments
\newtheorem{theorem}{Theorem}[section]
\newtheorem{proposition}[theorem]{Proposition}
\newtheorem{lemma}[theorem]{Lemma}
\newtheorem{corollary}[theorem]{Corollary}
\theoremstyle{definition}
\newtheorem{definition}[theorem]{Definition}
\newtheorem{example}[theorem]{Example}
\theoremstyle{remark}
\newtheorem{remark}[theorem]{Remark}

% Shortcuts
\newcommand{\C}{\mathbb{C}}
\newcommand{\Z}{\mathbb{Z}}
\newcommand{\N}{\mathbb{N}}
\newcommand{\gl}{\mathfrak{gl}}
\newcommand{\End}{\operatorname{End}}
\newcommand{\Hom}{\operatorname{Hom}}
\newcommand{\Sym}{\operatorname{Sym}}
\newcommand{\Span}{\operatorname{span}}
\newcommand{\sgn}{\operatorname{sgn}}
\newcommand{\SYT}{\operatorname{SYT}}
\newcommand{\SSYT}{\operatorname{SSYT}}
\newcommand{\shape}{\operatorname{shape}}
\newcommand{\cont}{\operatorname{cont}}
\newcommand{\ax}{\operatorname{ax}}
\newcommand{\Rep}{\mathcal{R}\text{ep}}

\title{Classical Schur--Weyl Duality\\with Explicit Matrix Elements}
\author{Notes}
\date{}

\begin{document}
\maketitle

\begin{abstract}
We give a self-contained account of classical Schur--Weyl duality for
$\gl_N$ and $S_n$, emphasising the explicit construction of matrix elements
for both actions. The Robinson--Schensted--Knuth correspondence emerges
as the combinatorial incarnation of the decomposition: the pair $(P, Q)$
produced by RSK literally \emph{is} the Schur--Weyl decomposition, giving
coordinates in the $\gl_N$ and $S_n$ factors simultaneously.
\end{abstract}

\tableofcontents

\part{Classical Schur--Weyl Duality}

%% ====================================================================
\section{The two commuting actions}
\label{sec:setup}
%% ====================================================================

Fix a positive integer $N$ and let $V = \C^N$ with standard basis
$e_1, \ldots, e_N$. For any positive integer $n$, the tensor power
$V^{\otimes n}$ has dimension $N^n$ with basis
\[
  e_\mathbf{i} \;=\; e_{i_1} \otimes e_{i_2} \otimes \cdots \otimes e_{i_n},
  \qquad \mathbf{i} = (i_1, \ldots, i_n) \in \{1,\ldots,N\}^n.
\]
We call $\mathbf{i}$ a \emph{word} of length $n$ in the alphabet
$\{1, \ldots, N\}$.

Two groups act naturally on $V^{\otimes n}$.

\subsection{The $\mathrm{GL}_N$ action}

The general linear group $\mathrm{GL}_N(\C)$ acts \emph{diagonally}:
\[
  g \cdot (e_{i_1} \otimes \cdots \otimes e_{i_n})
  = g\,e_{i_1} \otimes \cdots \otimes g\,e_{i_n}.
\]
In matrix elements, if $g = (g_{ab})$ then
\[
  g \cdot e_\mathbf{i}
  = \sum_{\mathbf{j} \in \{1,\ldots,N\}^n}
    \Bigl(\prod_{k=1}^{n} g_{j_k, i_k}\Bigr)\, e_\mathbf{j}.
\]
The corresponding Lie algebra action of $\gl_N$ is given by the
matrix units $E_{ab}$ ($1 \leqslant a, b \leqslant N$) acting as
\begin{equation}\label{eq:glN-action}
  E_{ab} \cdot e_\mathbf{i}
  = \sum_{k=1}^{n} \delta_{b, i_k}\,
    e_{i_1} \otimes \cdots \otimes e_{i_{k-1}}
    \otimes e_a
    \otimes e_{i_{k+1}} \otimes \cdots \otimes e_{i_n}.
\end{equation}
That is, $E_{ab}$ scans each tensor factor and replaces $e_b$ by $e_a$
wherever it finds it.

\subsection{The $S_n$ action}

The symmetric group $S_n$ acts by \emph{permuting tensor positions}:
\begin{equation}\label{eq:Sn-action}
  \sigma \cdot (e_{i_1} \otimes \cdots \otimes e_{i_n})
  = e_{i_{\sigma^{-1}(1)}} \otimes \cdots \otimes e_{i_{\sigma^{-1}(n)}}.
\end{equation}
Equivalently, $\sigma$ sends the word $\mathbf{i} = (i_1, \ldots, i_n)$
to the word $\sigma \cdot \mathbf{i} = (i_{\sigma^{-1}(1)}, \ldots, i_{\sigma^{-1}(n)})$.
Since $S_n$ is generated by adjacent transpositions $s_k = (k,\; k{+}1)$
for $k = 1, \ldots, n{-}1$, it suffices to specify:
\[
  s_k \cdot e_\mathbf{i}
  = e_{i_1} \otimes \cdots \otimes e_{i_{k+1}} \otimes e_{i_k}
    \otimes \cdots \otimes e_{i_n},
\]
which simply swaps positions $k$ and $k{+}1$.

\subsection{Commutativity}

\begin{proposition}\label{prop:commute}
The actions of $\mathrm{GL}_N$ and $S_n$ on $V^{\otimes n}$ commute.
\end{proposition}

\begin{proof}
For any $g \in \mathrm{GL}_N$ and $\sigma \in S_n$,
\begin{align*}
  g \cdot (\sigma \cdot e_\mathbf{i})
  &= g \cdot (e_{i_{\sigma^{-1}(1)}} \otimes \cdots
       \otimes e_{i_{\sigma^{-1}(n)}}) \\
  &= g\,e_{i_{\sigma^{-1}(1)}} \otimes \cdots
       \otimes g\,e_{i_{\sigma^{-1}(n)}},
\end{align*}
while
\begin{align*}
  \sigma \cdot (g \cdot e_\mathbf{i})
  &= \sigma \cdot (g\,e_{i_1} \otimes \cdots \otimes g\,e_{i_n}) \\
  &= g\,e_{i_{\sigma^{-1}(1)}} \otimes \cdots
       \otimes g\,e_{i_{\sigma^{-1}(n)}}.
\end{align*}
The key: $g$ acts within each tensor factor (it doesn't see position),
while $\sigma$ rearranges factors (it doesn't see which basis vector
is in each slot).
\end{proof}

%% ====================================================================
\section{The decomposition}
\label{sec:decomposition}
%% ====================================================================

\subsection{Partitions and Young diagrams}

A \emph{partition} $\lambda$ of $n$, written $\lambda \vdash n$, is a
weakly decreasing sequence $\lambda = (\lambda_1 \geqslant \lambda_2
\geqslant \cdots \geqslant \lambda_\ell > 0)$ with
$|\lambda| = \lambda_1 + \cdots + \lambda_\ell = n$. We write
$\ell(\lambda) = \ell$ for the number of parts. The \emph{Young diagram}
of $\lambda$ consists of $\lambda_i$ left-justified boxes in row $i$
(using English convention, with row 1 at the top).

\begin{example}
The partition $\lambda = (3,2,1) \vdash 6$ has Young diagram
\[
  \ydiagram{3,2,1}
\]
\end{example}

\subsection{Standard and semistandard tableaux}

\begin{definition}
A \emph{standard Young tableau} (SYT) of shape $\lambda \vdash n$ is a
filling of the Young diagram of $\lambda$ with the numbers $1, 2,
\ldots, n$, each used exactly once, such that entries increase along
each row (left to right) and down each column (top to bottom).
\end{definition}

\begin{definition}
A \emph{semistandard Young tableau} (SSYT) of shape $\lambda$ with
entries in $\{1, \ldots, N\}$ is a filling of the Young diagram of
$\lambda$ such that entries weakly increase along each row and
strictly increase down each column.
\end{definition}

We write $\SYT(\lambda)$ for the set of standard tableaux of shape
$\lambda$, and $\SSYT(\lambda, N)$ for the set of semistandard tableaux
of shape $\lambda$ with entries in $\{1, \ldots, N\}$.

\begin{example}\label{ex:tableaux-21}
For $\lambda = (2,1)$ and $N = 2$:
\[
  \SYT(2,1) = \Bigl\{\;
    \ytableaushort{12,3}\;,\quad
    \ytableaushort{13,2}\;
  \Bigr\},
  \qquad
  \SSYT\bigl((2,1),\, 2\bigr) = \Bigl\{\;
    \ytableaushort{11,2}\;,\quad
    \ytableaushort{12,2}\;
  \Bigr\}.
\]
\end{example}

The dimensions of the irreducible representations are given by:
\begin{itemize}
\item $\dim S^\lambda = |\SYT(\lambda)| = \dfrac{n!}{\prod_{u \in \lambda} h(u)}$
  \quad (hook length formula),
\item $\dim W^\lambda = |\SSYT(\lambda, N)|
  = \prod_{1 \leqslant i < j \leqslant N}
    \dfrac{\lambda_i - \lambda_j + j - i}{j - i}$
  \quad (Weyl dimension formula),
\end{itemize}
where $h(u)$ is the hook length of the box $u$ in the Young diagram
of~$\lambda$.

\subsection{The Schur--Weyl decomposition}

\begin{theorem}[Schur--Weyl duality]\label{thm:SW}
As a $(\mathrm{GL}_N, S_n)$-bimodule,
\begin{equation}\label{eq:SW}
  V^{\otimes n}
  \;\cong\;
  \bigoplus_{\substack{\lambda \vdash n \\ \ell(\lambda) \leqslant N}}
  W^\lambda \otimes S^\lambda,
\end{equation}
where:
\begin{itemize}
\item $W^\lambda$ is the irreducible polynomial representation of
  $\mathrm{GL}_N$ of highest weight $\lambda$ (the \emph{Weyl module}
  or \emph{Schur module}),
\item $S^\lambda$ is the irreducible representation of $S_n$ indexed
  by $\lambda$ (the \emph{Specht module}).
\end{itemize}
Moreover:
\begin{enumerate}[label=(\alph*)]
\item The natural map
  $\C[S_n] \to \End_{\gl_N}(V^{\otimes n})$
  is surjective (and injective when $n \leqslant N$).
\item The natural map
  $U(\gl_N) \to \End_{S_n}(V^{\otimes n})$
  is surjective.
\end{enumerate}
\end{theorem}

\begin{remark}
Checking dimensions: on the left, $\dim V^{\otimes n} = N^n$. On the
right,
\[
  \sum_{\substack{\lambda \vdash n \\ \ell(\lambda) \leqslant N}}
  |\SSYT(\lambda, N)| \cdot |\SYT(\lambda)| = N^n.
\]
This identity can be proved bijectively via the RSK correspondence
(Section~\ref{sec:RSK}).
\end{remark}

%% ====================================================================
\section{Explicit matrix elements for $S_n$: Young's orthogonal form}
\label{sec:Young}
%% ====================================================================

We now describe how $S_n$ acts on each Specht module $S^\lambda$ in
an explicit basis indexed by standard Young tableaux.

\subsection{Content and axial distance}

\begin{definition}\label{def:content}
Let $T$ be a standard Young tableau. The \emph{content} of the box in
row $r$ and column $c$ is $\cont(r,c) = c - r$. For each entry
$k \in \{1, \ldots, n\}$, we write $\cont_T(k)$ for the content of
the box containing $k$ in $T$.

The \emph{axial distance} from $k$ to $k{+}1$ in $T$ is
\begin{equation}\label{eq:axial}
  \rho_T(k) = \cont_T(k{+}1) - \cont_T(k).
\end{equation}
\end{definition}

\begin{example}
In the standard tableau $T = \;\ytableaushort{13,2}\;$ of shape
$(2,1)$: the entry $1$ is in box $(1,1)$ with content $0$, the entry $2$ is
in box $(2,1)$ with content $-1$, and $3$ is in box $(1,2)$ with
content $1$. So $\rho_T(1) = \cont_T(2) - \cont_T(1) = -1 - 0 = -1$
and $\rho_T(2) = \cont_T(3) - \cont_T(2) = 1 - (-1) = 2$.
\end{example}

\begin{remark}
The axial distance $\rho_T(k)$ is always a nonzero integer. It is
positive when $k{+}1$ appears strictly to the right and weakly below
$k$ (more precisely, when the content increases), and negative
otherwise. The case $|\rho_T(k)| = 1$ occurs exactly when $k$ and
$k{+}1$ are in the same row or the same column; swapping them would
violate the standard tableau condition, so $s_k \cdot T$ is undefined
in that case.
\end{remark}

\subsection{The orthogonal representation}

\begin{definition}\label{def:Young-orth}
For each $\lambda \vdash n$, Young's \emph{orthogonal representation}
is the representation of $S_n$ on the vector space
$\widetilde{S}^\lambda = \bigoplus_{T \in \SYT(\lambda)} \C \cdot v_T$
defined by the action of the adjacent transpositions:
\begin{equation}\label{eq:Young-orth}
  s_k \cdot v_T = \frac{1}{\rho_T(k)}\, v_T
  + \sqrt{1 - \frac{1}{\rho_T(k)^2}}\; v_{s_k \cdot T},
\end{equation}
where $s_k \cdot T$ is the tableau obtained from $T$ by swapping the
entries $k$ and $k{+}1$.

If $s_k \cdot T$ is not a valid standard tableau (i.e.,
$|\rho_T(k)| = 1$), then the second term is absent and
$s_k \cdot v_T = \pm v_T$.
\end{definition}

\begin{remark}
There is also Young's \emph{seminormal form} which uses the same
eigenvectors but different normalization:
\[
  s_k \cdot v_T^{\mathrm{sn}} = \frac{1}{\rho_T(k)}\, v_T^{\mathrm{sn}}
  + \Bigl(1 + \frac{1}{\rho_T(k)}\Bigr) v_{s_k \cdot T}^{\mathrm{sn}}.
\]
The orthogonal form has the advantage that the representing matrices
are orthogonal.
\end{remark}

\begin{example}\label{ex:S3-21}
For $\lambda = (2,1) \vdash 3$, we have two standard tableaux:
\[
  T_1 = \;\ytableaushort{12,3}\;, \qquad
  T_2 = \;\ytableaushort{13,2}\;.
\]
The axial distances are $\rho_{T_1}(1) = 1$ (so $s_1 \cdot T_1 =
T_1$ is forced), $\rho_{T_1}(2) = -2$, $\rho_{T_2}(1) = -1$ (so
$s_1 \cdot T_2 = -T_2$ is forced), and $\rho_{T_2}(2) = 2$. Thus in
the basis $(v_{T_1}, v_{T_2})$:
\[
  s_1 = \begin{pmatrix} 1 & 0 \\ 0 & -1 \end{pmatrix},
  \qquad
  s_2 = \begin{pmatrix} -\frac{1}{2} & \frac{\sqrt{3}}{2} \\[4pt]
    \frac{\sqrt{3}}{2} & \frac{1}{2} \end{pmatrix}.
\]
One can verify that $s_1^2 = s_2^2 = I$ and $(s_1 s_2)^3 = I$,
confirming we have a representation of $S_3$.
\end{example}

%% ====================================================================
\section{Explicit matrix elements for $\gl_N$: the Gel'fand--Tsetlin basis}
\label{sec:GT}
%% ====================================================================

\subsection{Gel'fand--Tsetlin patterns}

The branching rule for $\gl_N \supset \gl_{N-1}$ says that the
irreducible $\gl_N$-module $W^\lambda$ with highest weight
$\lambda = (\lambda_1, \ldots, \lambda_N)$ decomposes as
\[
  W^\lambda\big|_{\gl_{N-1}}
  \;=\; \bigoplus_\mu W^\mu,
\]
where $\mu = (\mu_1, \ldots, \mu_{N-1})$ runs over all partitions
satisfying the \emph{interlacing condition}
\[
  \lambda_1 \geqslant \mu_1 \geqslant \lambda_2 \geqslant \mu_2
  \geqslant \cdots \geqslant \mu_{N-1} \geqslant \lambda_N.
\]
Iterating this branching down the chain
$\gl_N \supset \gl_{N-1} \supset \cdots \supset \gl_1$, each basis
vector is indexed by a triangular array:

\begin{definition}
A \emph{Gel'fand--Tsetlin (GT) pattern} of shape $\lambda$ (with
$N$ rows) is a triangular array
\[
  \Lambda = \begin{pmatrix}
    \lambda_{N,1} & & \lambda_{N,2} & & \cdots & & \lambda_{N,N} \\
    & \lambda_{N-1,1} & & \lambda_{N-1,2} & & \cdots & \lambda_{N-1,N-1} \\
    & & \ddots & & & & \\
    & & & \lambda_{1,1} & & &
  \end{pmatrix}
\]
with $(\lambda_{N,1}, \ldots, \lambda_{N,N}) = \lambda$ and satisfying
the interlacing conditions
$\lambda_{k,i} \geqslant \lambda_{k-1,i} \geqslant \lambda_{k,i+1}$
for all valid indices.
\end{definition}

\begin{remark}
The set of GT patterns of shape $\lambda$ is in natural bijection
with $\SSYT(\lambda, N)$: given a semistandard tableau $T$, the $k$-th
row of the GT pattern records the shape of the subtableau of $T$
consisting of entries $\leqslant k$.
\end{remark}

\subsection{The $\gl_N$ action in the GT basis}

Write $v_\Lambda$ for the basis vector of $W^\lambda$ indexed by the
GT pattern $\Lambda$. The Lie algebra $\gl_N$ is generated by the
matrix units $E_{k,k}$ (Cartan subalgebra) and $E_{k,k+1}$,
$E_{k+1,k}$ (raising and lowering operators).

\begin{proposition}[Gel'fand--Tsetlin formulae]
\label{prop:GT-action}
The action of $\gl_N$ on the GT basis is:
\begin{align}
  E_{k,k}\, v_\Lambda
  &= \biggl(\sum_{i=1}^k \lambda_{k,i} - \sum_{i=1}^{k-1}
     \lambda_{k-1,i}\biggr)\, v_\Lambda,
  \label{eq:GT-diag} \\[6pt]
  E_{k,k+1}\, v_\Lambda
  &= -\sum_{i=1}^k
    \frac{\prod_{j=1}^{k+1}(\lambda_{k,i} - \lambda_{k+1,j} + j - i)
    \;\prod_{j=1}^{k-1}(\lambda_{k,i} - \lambda_{k-1,j} + j - i)}
    {\prod_{j \neq i}^{k}(\lambda_{k,i} - \lambda_{k,j} + j - i)}
    \; v_{\Lambda + \delta_{k,i}},
  \label{eq:GT-raise} \\[6pt]
  E_{k+1,k}\, v_\Lambda
  &= \sum_{i=1}^k
    \frac{\prod_{j=1}^{k+1}(\lambda_{k,i} - \lambda_{k+1,j} + j - i)
    \;\prod_{j=1}^{k-1}(\lambda_{k,i} - \lambda_{k-1,j} + j - i)}
    {\prod_{j \neq i}^{k}(\lambda_{k,i} - \lambda_{k,j} + j - i)}
    \; v_{\Lambda - \delta_{k,i}},
  \label{eq:GT-lower}
\end{align}
where $\Lambda \pm \delta_{k,i}$ denotes the GT pattern obtained by
replacing $\lambda_{k,i}$ with $\lambda_{k,i} \pm 1$ (and the term
is zero if the result is not a valid GT pattern).
\end{proposition}

The eigenvalue in~\eqref{eq:GT-diag} is the weight: $E_{k,k}$ counts
how many entries equal to $k$ appear in the corresponding semistandard
tableau.

\begin{remark}[Shifted coordinates]\label{rem:shifted}
The GT formulae simplify in \emph{shifted coordinates}
$\ell_{k,i} = \lambda_{k,i} - i$, which is the content of the
rightmost box in row~$i$ of the partition
$\lambda^{(k)} = (\lambda_{k,1}, \ldots, \lambda_{k,k})$.
The interlacing condition becomes
$\ell_{k+1,j} \geqslant \ell_{k,j} \geqslant \ell_{k+1,j+1}$,
and every factor in the raising and lowering formulae becomes a
difference of shifted entries:
\begin{align}
  E_{k,k+1}\, v_\Lambda
  &= -\sum_{i=1}^k
    \frac{\displaystyle\prod_{j=1}^{k+1}(\ell_{k,i} - \ell_{k+1,j})
    \;\prod_{j=1}^{k-1}(\ell_{k,i} - \ell_{k-1,j})}
    {\displaystyle\prod_{\substack{j=1 \\ j \neq i}}^{k}
    (\ell_{k,i} - \ell_{k,j})}
    \; v_{\Lambda + \delta_{k,i}},
  \label{eq:GT-raise-shifted} \\[6pt]
  E_{k+1,k}\, v_\Lambda
  &= \sum_{i=1}^k
    \frac{\displaystyle\prod_{j=1}^{k+1}(\ell_{k,i} - \ell_{k+1,j})
    \;\prod_{j=1}^{k-1}(\ell_{k,i} - \ell_{k-1,j})}
    {\displaystyle\prod_{\substack{j=1 \\ j \neq i}}^{k}
    (\ell_{k,i} - \ell_{k,j})}
    \; v_{\Lambda - \delta_{k,i}}.
  \label{eq:GT-lower-shifted}
\end{align}
The structure of each coefficient is now transparent:
\begin{itemize}
\item \textbf{Numerator}: product of distances from $\ell_{k,i}$ to
  the shifted entries on the rows immediately above
  ($\ell_{k+1,1}, \ldots, \ell_{k+1,k+1}$) and below
  ($\ell_{k-1,1}, \ldots, \ell_{k-1,k-1}$) in the GT pattern.
\item \textbf{Denominator}: Vandermonde-type product of distances
  between $\ell_{k,i}$ and all other shifted entries on the same
  row~$k$.
\end{itemize}
The Weyl dimension formula takes the same shape:
$\dim W^\lambda = \prod_{1 \leqslant i < j \leqslant N}
\frac{\ell_i - \ell_j}{j - i}$, where $\ell_i = \lambda_i - i$.
The GT coefficients are \emph{local} versions of this formula:
where the dimension involves all pairs of shifted entries across
the full partition, each GT coefficient involves only the entry
being moved and its neighbours in adjacent rows.
\end{remark}

%% ====================================================================
\section{The Robinson--Schensted--Knuth correspondence}
\label{sec:RSK}
%% ====================================================================

We now describe the bijection that ties everything together.

\subsection{Row insertion}

Given a semistandard tableau $P$ and a positive integer $x$, the
\emph{row insertion} $P \leftarrow x$ produces a new semistandard
tableau by the following algorithm:

\begin{enumerate}
\item Start with the first row of $P$.
\item If $x$ is weakly larger than all entries in the current row,
  place $x$ at the end of the row. Terminate.
\item Otherwise, find the leftmost entry $y$ in the current row with
  $y > x$. Replace $y$ by $x$ (``$x$ bumps $y$'').
\item Repeat from step 2 with $y$ and the next row down.
\end{enumerate}

The key property is that insertion preserves the semistandard
condition.

\subsection{The RSK correspondence}

\begin{definition}
The \emph{Robinson--Schensted--Knuth (RSK) correspondence} is a
bijection
\[
  \mathrm{RSK} \;:\; \{1, \ldots, N\}^n
  \;\xrightarrow{\;\sim\;}\;
  \bigsqcup_{\substack{\lambda \vdash n \\ \ell(\lambda) \leqslant N}}
  \SSYT(\lambda, N) \times \SYT(\lambda)
\]
defined as follows. Given a word $\mathbf{i} = (i_1, \ldots, i_n)$,
construct a pair of tableaux $(P, Q)$ by iterating:
\begin{itemize}
\item Start with $P_0 = Q_0 = \varnothing$ (empty tableaux).
\item At step $k = 1, \ldots, n$: insert $i_k$ into $P_{k-1}$ to get
  $P_k = P_{k-1} \leftarrow i_k$. This insertion adds a new box to
  the shape; place $k$ in that box to form $Q_k$.
\end{itemize}
Set $P = P_n$ (the \emph{insertion tableau}) and $Q = Q_n$ (the
\emph{recording tableau}).
\end{definition}

\begin{example}\label{ex:RSK}
Let $N = 2$, $n = 3$, and $\mathbf{i} = (2, 1, 1)$.
\begin{itemize}
\item Insert $2$: $P_1 = \;\ytableaushort{2}\;$, \;
  $Q_1 = \;\ytableaushort{1}\;$.
\item Insert $1$: $1$ bumps $2$ from row 1, $2$ goes to row 2.
  $P_2 = \;\ytableaushort{1,2}\;$, \;
  $Q_2 = \;\ytableaushort{1,2}\;$.
\item Insert $1$: we bump the leftmost entry \emph{strictly} greater
  than $1$ in row~1. There is no such entry (the only entry is $1$
  itself), so $1$ appends to the end of row~1.
  $P_3 = \;\ytableaushort{11,2}\;$, \;
  $Q_3 = \;\ytableaushort{13,2}\;$.
\end{itemize}
So $\mathrm{RSK}(2,1,1) = (P, Q)$ with $P = \;\ytableaushort{11,2}\;$ and
$Q = \;\ytableaushort{13,2}\;$, both of shape $\lambda = (2,1)$.
\end{example}

\subsection{Properties of RSK}

\begin{proposition}\label{prop:RSK-props}
The RSK correspondence has the following properties:
\begin{enumerate}[label=(\alph*)]
\item $P$ is a semistandard Young tableau with entries in $\{1,\ldots,N\}$.
\item $Q$ is a standard Young tableau.
\item $\shape(P) = \shape(Q) = \lambda$ for some partition $\lambda
  \vdash n$ with $\ell(\lambda) \leqslant N$.
\item The correspondence is a bijection.
\item The \emph{weight} of $P$ (the tuple recording how many times
  each letter appears) equals the weight of the word $\mathbf{i}$;
  specifically, the number of entries equal to $a$ in $P$ equals the
  number of occurrences of $a$ in $\mathbf{i}$.
\end{enumerate}
\end{proposition}

%% ====================================================================
\section{RSK is the Schur--Weyl decomposition}
\label{sec:RSK-is-SW}
%% ====================================================================

We now state the main interpretive result: the RSK correspondence is
not merely a combinatorial bijection, but the \emph{explicit
realisation} of the Schur--Weyl decomposition~\eqref{eq:SW}.

\subsection{The decomposition, concretely}

The tensor product basis element $e_\mathbf{i} \in V^{\otimes n}$
belongs to the summand $W^\lambda \otimes S^\lambda$
in~\eqref{eq:SW}, where $\lambda = \shape(P) = \shape(Q)$ and
$(P, Q) = \mathrm{RSK}(\mathbf{i})$. More precisely:

\begin{theorem}[RSK as Schur--Weyl]\label{thm:RSK-SW}
Under the isomorphism~\eqref{eq:SW}, the basis element $e_\mathbf{i}$
corresponds to
\[
  e_\mathbf{i} \;\longleftrightarrow\;
  w_P \otimes v_Q
  \;\in\; W^\lambda \otimes S^\lambda,
\]
where:
\begin{itemize}
\item $\lambda = \shape(\mathrm{RSK}(\mathbf{i}))$ determines which
  summand,
\item $P = P(\mathbf{i})$, the insertion tableau, gives coordinates
  in the $\mathrm{GL}_N$-irrep $W^\lambda$ (via the
  GT-basis/SSYT-basis identification),
\item $Q = Q(\mathbf{i})$, the recording tableau, gives coordinates
  in the $S_n$-irrep $S^\lambda$ (via the SYT basis).
\end{itemize}
\end{theorem}

\begin{remark}\label{rem:not-isometry}
The correspondence is not an isometry with respect to the naive inner
products---the tensor product basis is orthonormal while the RSK
correspondence identifies it with a \emph{different} basis of
$\bigoplus_\lambda W^\lambda \otimes S^\lambda$. The precise
relationship involves the transition matrix between the tensor product
basis and the Schur--Weyl adapted basis, which can be computed
recursively from the RSK insertion algorithm.
\end{remark}

\subsection{What each part of RSK ``sees''}

The beauty of this interpretation is that the two tableaux capture
complementary information:

\begin{center}
\renewcommand{\arraystretch}{1.4}
\begin{tabular}{c|c|c}
  & \textbf{Insertion tableau $P$}
  & \textbf{Recording tableau $Q$} \\
\hline
\textbf{Type} & Semistandard (SSYT) & Standard (SYT) \\
\textbf{Entries} & Letters from $\{1,\ldots,N\}$ & Positions $\{1,\ldots,n\}$ \\
\textbf{Indexes} & $\mathrm{GL}_N$-irrep basis & $S_n$-irrep basis \\
\textbf{Sensitive to} & What letters appear & In what order they arrive \\
\textbf{Dimension count} & $|\SSYT(\lambda,N)| = \dim W^\lambda$
  & $|\SYT(\lambda)| = \dim S^\lambda$ \\
\textbf{Symmetry} & Changed by $\mathrm{GL}_N$ action
  & Changed by $S_n$ action \\
\end{tabular}
\end{center}

\subsection{How the two actions transform $(P, Q)$}

\begin{proposition}\label{prop:actions-on-PQ}
Under the RSK identification:
\begin{enumerate}[label=(\alph*)]
\item The $S_n$-action permutes the recording tableau $Q$ (changing
  the $S_n$ coordinates) while preserving the insertion tableau $P$
  (the $\mathrm{GL}_N$ coordinates are invisible to $S_n$).
\item The $\mathrm{GL}_N$-action transforms the insertion tableau $P$
  (changing the $\mathrm{GL}_N$ coordinates) while preserving the
  recording tableau $Q$ (the $S_n$ coordinates are invisible to
  $\mathrm{GL}_N$).
\end{enumerate}
This is exactly the content of Schur--Weyl duality: the two actions
commute because they act on complementary factors of $W^\lambda
\otimes S^\lambda$.
\end{proposition}

\begin{remark}
Statement (a) requires care: $S_n$ does not simply permute entries of
$Q$ in a naive sense. Rather, the RSK correspondence intertwines the
$S_n$-action on words with a well-defined action on standard tableaux
that corresponds to the representation of $S_n$ on $S^\lambda$.
The precise statement involves the \emph{Knuth equivalence}: two words
$\mathbf{i}$ and $\mathbf{j}$ have the same insertion tableau $P$ if
and only if they are Knuth-equivalent.
\end{remark}

\subsection{The dimension identity}

The bijectivity of RSK immediately gives the counting identity:
\begin{equation}\label{eq:dim-count}
  N^n
  = |\{1,\ldots,N\}^n|
  = \sum_{\substack{\lambda \vdash n \\ \ell(\lambda) \leqslant N}}
    |\SSYT(\lambda, N)| \cdot |\SYT(\lambda)|
  = \sum_{\substack{\lambda \vdash n \\ \ell(\lambda) \leqslant N}}
    \dim W^\lambda \cdot \dim S^\lambda,
\end{equation}
which is the dimension check for the Schur--Weyl
decomposition~\eqref{eq:SW}.

\begin{example}
For $N = 2$, $n = 3$: we need partitions of $3$ with at most $2$
parts, namely $\lambda = (3)$ and $\lambda = (2,1)$.
\begin{itemize}
\item $\lambda = (3)$: $|\SSYT((3), 2)| = |\{111, 112, 122, 222\}| = 4$,
  $|\SYT((3))| = 1$. Contribution: $4 \cdot 1 = 4$.
\item $\lambda = (2,1)$: $|\SSYT((2,1), 2)| = 2$ (see
  Example~\ref{ex:tableaux-21}), $|\SYT((2,1))| = 2$.
  Contribution: $2 \cdot 2 = 4$.
\end{itemize}
Total: $4 + 4 = 8 = 2^3$. \checkmark
\end{example}

%% ====================================================================
\section{The bridge: Jucys--Murphy elements and spectral projectors}
\label{sec:bridge}
%% ====================================================================

The preceding sections construct the $S_n$ and $\gl_N$ matrix elements
\emph{intrinsically}: Young's orthogonal form
(Section~\ref{sec:Young}) uses the combinatorics of standard tableaux,
and the GT formulae (Section~\ref{sec:GT}) use the branching chain
$\gl_N \supset \gl_{N-1} \supset \cdots \supset \gl_1$.
Section~\ref{sec:RSK-is-SW} claims RSK provides the change of basis,
but never makes the embedding explicit. How do these abstract
representations actually sit inside $V^{\otimes n}$?

The missing ingredient is the \emph{Jucys--Murphy elements}, which
provide concrete projection operators from $V^{\otimes n}$ onto each
irreducible summand.

\subsection{Motivation: building a commutative subalgebra}

How would one construct a large commutative subalgebra of $\C[S_n]$?

\medskip\noindent\textbf{Strategy: use the chain
$S_1 \subset S_2 \subset \cdots \subset S_n$.}
The centre $Z(\C[S_k])$ commutes with \emph{everything} in $\C[S_k]$,
and in particular with $Z(\C[S_j])$ for $j \leqslant k$ (since
$\C[S_j] \subset \C[S_k]$). So the algebra generated by all the
centres $Z(\C[S_1]), Z(\C[S_2]), \ldots, Z(\C[S_n])$ is commutative.
This is the \emph{Gelfand--Tsetlin (GZ) subalgebra}.

Each centre $Z(\C[S_k])$ is spanned by class sums.
At step $k$, the ``new'' class sum not already in $\C[S_{k-1}]$ is
$X_k = \sum_{1 \leqslant i < j \leqslant k} (i\;j)$
(the sum of all transpositions in $S_k$), which is central in
$\C[S_k]$. Set $J_k = X_k - X_{k-1}$: this peels off the
contribution of the new element $k$, and is the sum of all
transpositions that \emph{involve} $k$.

The element $J_k$ measures how $k$ interacts with the earlier elements
$\{1, \ldots, k{-}1\}$. Since these interactions at different steps
$k$ and $l$ are mediated by disjoint ``new arrivals,'' one might hope
they commute---and indeed they do. More is true: the symmetric power
sums $J_1^r + \cdots + J_k^r$ are central in $\C[S_k]$ for every $r$,
so $J_1, \ldots, J_n$ actually generate the \emph{entire} GZ
subalgebra, not just a piece of it~\cite{OkounkovVershik}.

\subsection{The Jucys--Murphy elements}

\begin{definition}\label{def:JM}
The \emph{Jucys--Murphy elements} of $\C[S_n]$ are
\[
  J_1 = 0, \qquad
  J_k = (1\;k) + (2\;k) + \cdots + (k{-}1\;\; k),
  \quad k = 2, \ldots, n.
\]
Equivalently, $J_k = X_k - X_{k-1}$ where
$X_k = \sum_{1 \leqslant i < j \leqslant k} (i\;j) \in Z(\C[S_k])$.
\end{definition}

\begin{lemma}\label{lem:JM-commute}
The elements $J_1, \ldots, J_n$ commute pairwise.
\end{lemma}

\begin{proof}
Let $k < l$. We compute $[J_k, J_l] = \sum_{i < k}\sum_{j < l}
[(i\;k),\,(j\;l)]$ by sorting the terms of the double sum according
to whether the two transpositions share an index.

\emph{Case 1: $\{i,k\} \cap \{j,l\} = \varnothing$.}
The transpositions $(i\;k)$ and $(j\;l)$ are disjoint, hence commute.
Contribution: $0$.

\emph{Case 2: $i = j$ (so $i < k < l$).}
We get $[(i\;k),\,(i\;l)] = (i\;k\;l) - (i\;l\;k)$.

\emph{Case 3: $j = k$ (so $i < k < l$).}
We get $[(i\;k),\,(k\;l)] = (i\;l\;k) - (i\;k\;l)$.

No other overlaps occur: $i = l$ is impossible since $i < k < l$,
and $j = k$ with $j = i$ would force $i = k$.

For each value of $i < k$, the Case~2 and Case~3 contributions sum to
\[
  \bigl((i\;k\;l) - (i\;l\;k)\bigr)
  + \bigl((i\;l\;k) - (i\;k\;l)\bigr) = 0.
\]
So $[J_k, J_l] = 0$.
\end{proof}

The subalgebra of $\C[S_n]$ generated by $J_1, \ldots, J_n$ is called
the \emph{Gelfand--Tsetlin (GZ) subalgebra} for the chain
$S_1 \subset S_2 \subset \cdots \subset S_n$
\cite{OkounkovVershik}. By ``commutative subalgebra'' we mean that
the $J_k$ commute \emph{with each other}---they certainly do not
commute with arbitrary elements of $\C[S_n]$.

\begin{lemma}\label{lem:JM-recursion}
For $k \geqslant 2$,
\begin{equation}\label{eq:JM-recursion}
  J_k = s_{k-1}\, J_{k-1}\, s_{k-1} + s_{k-1}.
\end{equation}
\end{lemma}

\begin{proof}
We show $s_{k-1}\, J_{k-1}\, s_{k-1} = J_k - s_{k-1}$.
Recall $J_{k-1} = \sum_{i=1}^{k-2} (i,\, k{-}1)$ and
$s_{k-1} = (k{-}1,\, k)$. Conjugating each summand:
\[
  s_{k-1}\,(i,\, k{-}1)\, s_{k-1}
  = (k{-}1,\, k)\,(i,\, k{-}1)\,(k{-}1,\, k)
  = (i,\, k),
  \qquad i < k{-}1,
\]
since conjugation replaces every occurrence of $k{-}1$ by $k$.
Summing over $i = 1, \ldots, k{-}2$:
\[
  s_{k-1}\, J_{k-1}\, s_{k-1}
  = \sum_{i=1}^{k-2} (i,\, k)
  = J_k - (k{-}1,\, k)
  = J_k - s_{k-1}. \qedhere
\]
\end{proof}

\begin{proposition}\label{prop:JM-eigenvalues}
In the orthogonal representation $\widetilde{S}^\lambda$
(Definition~\ref{def:Young-orth}), the Jucys--Murphy element $J_k$
acts diagonally in the SYT basis:
\[
  J_k \cdot v_T = \cont_T(k)\, v_T,
\]
where $\cont_T(k) = c - r$ is the content of the box in row~$r$,
column~$c$ that contains entry~$k$ in~$T$
(Definition~\ref{def:content}).
\end{proposition}

\begin{proof}
We use the recursion of Lemma~\ref{lem:JM-recursion} and induction
on $k$.

\emph{Base case ($k = 1$):} $J_1 = 0$, and entry $1$ is always in
box $(1,1)$ with content $0$.

\emph{Inductive step:} Assume $J_{k-1}\, v_T = \cont_T(k{-}1)\, v_T$
for all $T$. Set $\rho = \rho_T(k{-}1) = \cont_T(k) - \cont_T(k{-}1)$,
and write $\alpha = 1/\rho$ and
$\beta = \sqrt{1 - 1/\rho^2}$. If $|\rho| = 1$ then $\beta = 0$ and
$s_{k-1}\, v_T = \alpha\, v_T$; the computation below applies with
$\beta = 0$. Otherwise let $T' = s_{k-1} \cdot T$ (the tableau with
entries $k{-}1$ and $k$ swapped). Note that $\cont_{T'}(k{-}1) =
\cont_T(k)$ and $\rho_{T'}(k{-}1) = -\rho$, so from
Young's orthogonal form:
\[
  s_{k-1}\, v_T = \alpha\, v_T + \beta\, v_{T'},
  \qquad
  s_{k-1}\, v_{T'} = -\alpha\, v_{T'} + \beta\, v_T.
\]

Now compute $s_{k-1}\, J_{k-1}\, s_{k-1}\, v_T$ step by step.
Writing $c_1 = \cont_T(k{-}1)$ and $c_2 = \cont_T(k)$:
\begin{align*}
  s_{k-1}\, v_T &= \alpha\, v_T + \beta\, v_{T'}, \\
  J_{k-1}(\alpha\, v_T + \beta\, v_{T'})
    &= \alpha\, c_1\, v_T + \beta\, c_2\, v_{T'}
    \qquad\text{(by induction)}, \\
  s_{k-1}(\alpha\, c_1\, v_T + \beta\, c_2\, v_{T'})
    &= \alpha\, c_1\,(\alpha\, v_T + \beta\, v_{T'})
     + \beta\, c_2\,(\beta\, v_T - \alpha\, v_{T'}) \\
    &= \bigl(\alpha^2 c_1 + \beta^2 c_2\bigr)\, v_T
     + \alpha\beta\bigl(c_1 - c_2\bigr)\, v_{T'}.
\end{align*}
Since $c_1 - c_2 = -\rho$ and $\alpha = 1/\rho$:
\begin{itemize}
\item the $v_{T'}$ coefficient is
  $\alpha\beta(-\rho) = -\beta$,
\item the $v_T$ coefficient is
  $\alpha^2 c_1 + (1 - \alpha^2)\, c_2 = c_2 + \alpha^2(c_1 - c_2)
  = c_2 - 1/\rho$.
\end{itemize}
Adding $s_{k-1}\, v_T = (1/\rho)\, v_T + \beta\, v_{T'}$:
\[
  J_k\, v_T
  = \bigl(c_2 - \tfrac{1}{\rho} + \tfrac{1}{\rho}\bigr)\, v_T
  + (-\beta + \beta)\, v_{T'}
  = c_2\, v_T
  = \cont_T(k)\, v_T.
\]
The off-diagonal terms cancel exactly.
\end{proof}

That is, the SYT basis of Young's orthogonal form is precisely the
\emph{simultaneous eigenbasis} of all Jucys--Murphy elements. The
\emph{content vector} of a standard tableau~$T$ is
$\mathbf{c}(T) = (\cont_T(1), \ldots, \cont_T(n))$.

\begin{remark}[What simultaneous diagonalisation achieves]
\label{rem:simdiag}
We first recall a general fact from linear algebra.

\medskip\noindent\textbf{The single-operator case.}
Let $A$ be a diagonalisable operator on a finite-dimensional vector
space $V$ with distinct eigenvalues $c_1, \ldots, c_m$.
For each $i$, define
\begin{equation}\label{eq:lagrange-proj}
  P_i = \prod_{j \neq i} \frac{A - c_j}{c_i - c_j}.
\end{equation}
This is a polynomial in $A$ of degree $m - 1$. We claim $P_i$ is
the projector onto the $c_i$-eigenspace.

\emph{Proof.} Let $v$ be an eigenvector with $Av = c_k v$.
Substituting $A = c_k$ into~\eqref{eq:lagrange-proj}:
\[
  P_i\, v = \prod_{j \neq i} \frac{c_k - c_j}{c_i - c_j}\, v.
\]
If $k = i$, every factor is $(c_i - c_j)/(c_i - c_j) = 1$, so
$P_i\, v = v$. If $k \neq i$, the factor $j = k$ gives
$(c_k - c_k)/(c_i - c_k) = 0$, so $P_i\, v = 0$.
Since eigenvectors span $V$, this determines $P_i$ completely:
it acts as the identity on the $c_i$-eigenspace and as zero on
all other eigenspaces.

In particular: $P_i^2 = P_i$ (applying the projector twice is
the same as once), $P_i P_j = 0$ for $i \neq j$ (projecting onto
$c_i$ then $c_j$ kills everything), and $\sum_{i=1}^m P_i = I$
(every vector belongs to some eigenspace). Furthermore, $A$ itself
is recovered as $A = \sum_{i=1}^m c_i\, P_i$ (the \emph{spectral
decomposition}).

\medskip\noindent\textbf{Multiple commuting operators.}
If $A_1, \ldots, A_n$ are simultaneously diagonalisable, the
projector onto the joint eigenspace with eigenvalue tuple
$(c_1, \ldots, c_n)$ is the product of the individual Lagrange
projectors:
\[
  P_{(c_1,\ldots,c_n)}
  = \prod_{k=1}^n \;\prod_{\substack{c \neq c_k \\
    c \text{ an eigenvalue of } A_k}}
    \frac{A_k - c}{c_k - c}.
\]
This is a polynomial in $A_1, \ldots, A_n$; it acts as the identity
on the target eigenspace and as zero on all others, so
$P^2 = P$.

For the Jucys--Murphy elements, this general procedure has a
remarkable consequence. The operators $J_1, \ldots, J_n$ live in
$\C[S_n]$, so their Lagrange projectors are elements of $\C[S_n]$
too---they are the primitive idempotents $e_Q$ of
Definition~\ref{def:primitive-idempotent}. Each eigenvalue tuple
$\mathbf{c}(T)$ is the content vector of a standard tableau~$T$,
and the content vector carries two levels of information:
\begin{itemize}
\item The \emph{multiset} of contents
  $\{\cont_T(1), \ldots, \cont_T(n)\}$ determines the shape
  $\lambda$ (which irreducible representation).
\item The \emph{sequence} of contents
  $(\cont_T(1), \ldots, \cont_T(n))$ determines the specific
  tableau $T$ (which basis vector within that irreducible).
\end{itemize}
So the joint eigenspace decomposition is \emph{finer} than the
irreducible decomposition: each irreducible $S^\lambda$ of
dimension $d_\lambda = |\SYT(\lambda)|$ splits into $d_\lambda$
one-dimensional eigenspaces, one per standard tableau of shape
$\lambda$. Grouping eigenspaces by shape recovers the
irreducible decomposition.

Applied to \emph{any} representation $V$ of $S_n$: the
idempotents $e_Q \in \C[S_n]$ act on $V$, and the subspaces
$e_Q \cdot V$ decompose $V$ into pieces labelled by standard
tableaux. This is one diagonalisation, performed once inside the
algebra, valid for all representations simultaneously.
\end{remark}

\begin{lemma}\label{lem:content-separates}
Different standard tableaux have different content vectors: if $T$
and $T'$ are standard tableaux (of any shapes) with
$\mathbf{c}(T) = \mathbf{c}(T')$, then $T = T'$.
\end{lemma}

\begin{proof}
A standard tableau $T$ is built step by step: at step $k$, the entry
$k$ is placed in an addable corner of the shape $\lambda^{(k-1)}$
(the shape of entries $1, \ldots, k{-}1$) to produce
$\lambda^{(k)}$. The content $\cont_T(k)$ records
\emph{which} corner was chosen.

The key observation: the addable corners of any Young diagram have
distinct contents (each lies on a different diagonal $c - r$). So
knowing $\cont_T(k)$ and the shape $\lambda^{(k-1)}$ determines the
shape $\lambda^{(k)}$ uniquely.

Since $\lambda^{(0)} = \varnothing$ for every SYT, induction on $k$
shows that the content vector $(\cont_T(1), \ldots, \cont_T(n))$
determines the entire sequence of shapes
$\varnothing = \lambda^{(0)} \subset \lambda^{(1)} \subset \cdots
\subset \lambda^{(n)}$, which is the tableau $T$.
\end{proof}

\begin{theorem}\label{thm:GZ-maximal}
The GZ subalgebra is a maximal commutative subalgebra of $\C[S_n]$.
\end{theorem}

\begin{proof}
By Artin--Wedderburn, $\C[S_n] \cong \bigoplus_{\lambda \vdash n}
\mathrm{Mat}_{d_\lambda}(\C)$ where $d_\lambda = |\SYT(\lambda)|$.
A commutative subalgebra of $\mathrm{Mat}_d(\C)$ has dimension at most
$d$ (since commuting matrices can be simultaneously
upper-triangularised, and diagonal matrices in that basis span a
$d$-dimensional commutative subalgebra). So any commutative
subalgebra of $\C[S_n]$ has dimension at most
$\sum_{\lambda \vdash n} d_\lambda$.

By Proposition~\ref{prop:JM-eigenvalues}, the GZ subalgebra acts on
each $\widetilde{S}^\lambda$ by diagonal matrices in the SYT basis
$\{v_T : T \in \SYT(\lambda)\}$. By
Lemma~\ref{lem:content-separates}, different basis vectors have
different joint eigenvalues, so the GZ subalgebra contains \emph{all}
diagonal matrices in this basis (any diagonal matrix is a polynomial
in commuting operators with distinct joint eigenvalues---use Lagrange
interpolation). Thus the GZ subalgebra contains a
$d_\lambda$-dimensional subalgebra in each block, giving total
dimension at least $\sum_\lambda d_\lambda$.

Since the GZ subalgebra is commutative
(Lemma~\ref{lem:JM-commute}) and achieves the maximum possible
dimension for a commutative subalgebra, it is maximal.
\end{proof}

\subsection{Spectral projectors}

Since the $J_k$ are simultaneously diagonalisable, we can build
projection operators from their spectral decomposition.

\begin{definition}\label{def:primitive-idempotent}
For each $Q \in \SYT(\lambda)$, the \emph{primitive idempotent} is
\[
  e_Q \;=\; \prod_{k=2}^{n}\;\prod_{\substack{c \neq \cont_Q(k) \\
    c \text{ an eigenvalue of } J_k}}
    \frac{J_k - c}{\cont_Q(k) - c}
  \;\;\in\; \C[S_n].
\]
\end{definition}

These idempotents satisfy:
\begin{enumerate}[label=(\roman*)]
\item $e_Q^2 = e_Q$ and $e_Q \, e_{Q'} = 0$ for $Q \neq Q'$;
\item $\displaystyle\sum_{Q} e_Q = 1$, the sum over all SYT of all
  shapes;
\item $c_\lambda := \displaystyle\sum_{Q \in \SYT(\lambda)} e_Q$ is the
  central idempotent projecting onto the $\lambda$-isotypic component.
\end{enumerate}

\begin{remark}
Alternatively, $c_\lambda$ can be expressed via characters:
\[
  c_\lambda = \frac{\dim S^\lambda}{n!}
  \sum_{\sigma \in S_n} \chi^\lambda(\sigma)\, \sigma,
\]
but the Jucys--Murphy construction has the advantage of producing the
\emph{primitive} idempotents $e_Q$ directly, without requiring knowledge
of the full character table.
\end{remark}

\subsection{Carving out the Weyl modules}

Since $e_Q \in \C[S_n]$ and the $\gl_N$-action commutes with $S_n$
(Proposition~\ref{prop:commute}), the subspace
$e_Q \cdot V^{\otimes n} \subset V^{\otimes n}$
is a $\gl_N$-submodule. This is the bridge:

\begin{proposition}\label{prop:JM-bridge}
For each $Q \in \SYT(\lambda)$ with $\ell(\lambda) \leqslant N$:
\begin{enumerate}[label=(\alph*)]
\item $e_Q \cdot V^{\otimes n}$ is an irreducible $\gl_N$-module
  isomorphic to $W^\lambda$.
\item Under this isomorphism, the GT basis vector $v_\Lambda$
  corresponds to $e_Q \cdot e_\mathbf{i}$, where $\mathbf{i}$ is any
  word with $\mathrm{RSK}(\mathbf{i}) = (P, Q)$ and $P$ is the SSYT
  corresponding to $\Lambda$.
\item Varying $Q$ over $\SYT(\lambda)$ gives the $\dim S^\lambda$
  distinct copies of $W^\lambda$ constituting the $\lambda$-isotypic
  component of $V^{\otimes n}$.
\end{enumerate}
\end{proposition}

This resolves the gap identified in Remark~\ref{rem:not-isometry}: the
isomorphism~\eqref{eq:SW} is realised concretely by $e_\mathbf{i}
\mapsto \sum_Q e_Q \cdot e_\mathbf{i}$, decomposing each tensor into
its projections onto the copies of $W^\lambda$ labelled by standard
tableaux.

\subsection{How the two sets of matrix elements act on $V^{\otimes n}$}

The explicit formulae of Sections~\ref{sec:Young} and~\ref{sec:GT}
now have a concrete home:
\begin{itemize}
\item \textbf{The GT formulae} (Proposition~\ref{prop:GT-action})
  describe the $\gl_N$-action \emph{within} each projected subspace
  $e_Q \cdot V^{\otimes n}$, for a fixed~$Q$.
\item \textbf{Young's orthogonal form}
  (equation~\eqref{eq:Young-orth}) describes how the $S_n$-action
  \emph{rotates between} different copies: a transposition $s_k$
  sends $e_Q \cdot V^{\otimes n}$ to a linear combination of
  $e_Q \cdot V^{\otimes n}$ and $e_{s_k \cdot Q} \cdot V^{\otimes n}$,
  with coefficients given by~\eqref{eq:Young-orth}.
\end{itemize}

\subsection{Duality of branching chains}

The two sets of matrix elements arise from two dual branching chains
acting on orthogonal aspects of the tensor space:

\begin{center}
\renewcommand{\arraystretch}{1.4}
\begin{tabular}{c|c|c}
  & \textbf{$S_n$ chain} & \textbf{$\gl_N$ chain} \\
\hline
\textbf{Chain}
  & $S_1 \subset S_2 \subset \cdots \subset S_n$
  & $\gl_1 \subset \gl_2 \subset \cdots \subset \gl_N$ \\
\textbf{Acts on}
  & Positions $\{1, \ldots, n\}$
  & Values $\{1, \ldots, N\}$ \\
\textbf{On $V^{\otimes n}$}
  & Permutes which slot each value occupies
  & Mixes which value each slot contains \\
\textbf{Branching gives}
  & SYT basis (growth of shape by position)
  & GT basis (growth of shape by value) \\
\textbf{GZ generators}
  & Jucys--Murphy elements $J_k$
  & Cartan elements of $\gl_k$ \\
\textbf{Eigenvalues}
  & Contents $\cont_T(k)$
  & Weights $\textstyle\sum_i \lambda_{k,i} - \sum_i \lambda_{k-1,i}$
\end{tabular}
\end{center}

These two branching structures commute because positions and values
are independent: permuting tensor positions does not interfere with
mixing basis vectors within each position. The Schur--Weyl
decomposition is the \emph{simultaneous spectral decomposition} with
respect to both commuting families of operators.

%% ====================================================================
\section{Connection to the categorical framework}
\label{sec:categorical}
%% ====================================================================

Following Davydov--Molev~\cite{DM}, the classical Schur--Weyl duality
can be reformulated categorically as follows.

The sequence of group algebras $k[S_*] = \{k[S_n] \mid n \geqslant 0\}$
with the natural maps $\mu_{m,n}: k[S_m] \otimes k[S_n] \to k[S_{m+n}]$
induced by $S_m \times S_n \hookrightarrow S_{m+n}$ forms a
\emph{multiplicative sequence of algebras}. The associated
\emph{Schur--Weyl category} $\mathcal{S} = \mathcal{C}(k[S_*])$ is
a free symmetric monoidal category generated by one object, with
endomorphism algebras $\End_\mathcal{S}(X^{\otimes n}) = k[S_n]$.

The \emph{symmetric fiber functor} $F_N : \mathcal{S} \to
\mathcal{V}\!ect$ labelled by $N$ sends the generator $X$ to $V =
k^N$, and on morphisms gives the map $k[S_n] \to \End(V^{\otimes n})$.
This functor factors through the category of representations of
$\gl_N$ via the \emph{Schur--Weyl functor}
$SW_N : \mathcal{S} \to \Rep(\gl_N)$.

The classical Schur--Weyl duality (Theorem~\ref{thm:SW}) is precisely
the statement that $SW_N$ is \emph{full}: every $\gl_N$-equivariant
map between tensor powers of $V$ comes from the symmetric group.
The RSK correspondence provides the explicit combinatorial witness for
this fullness, by decomposing $V^{\otimes n}$ into irreducible
bimodule summands.

The power of the Davydov--Molev framework is that the same pattern
applies to quantum and affine deformations: replacing $k[S_*]$ by the
Hecke algebras $\mathrm{H}_*(q)$ gives the Hecke category whose fiber
functors factor through $U_q(\gl_N)$-modules; replacing by
degenerate affine Hecke algebras gives functors to Yangian modules;
and replacing by affine Hecke algebras gives functors to quantum
affine algebra modules. In each case, the universal property of the
Schur--Weyl category reduces the duality to a statement about
generators and relations.

%% ********************************************************************
\part{$q$-Deformed Schur--Weyl Duality}
\label{part:q}
%% ********************************************************************

Throughout Part~II, we fix a parameter $q \in \C^*$ with $q$ not a
root of unity. All constructions reduce to the classical case of
Part~I when $q \to 1$.

We use the notation
\[
  [x]_q = \frac{q^x - q^{-x}}{q - q^{-1}}
\]
for the \emph{$q$-number}. When $x$ is a positive integer $n$, this
is the \emph{$q$-integer} $[n]_q = q^{n-1} + q^{n-3} + \cdots +
q^{-(n-1)}$. Note $[x]_q \to x$ as $q \to 1$.

%% ====================================================================
\section{The two commuting actions ($q$-deformed)}
\label{sec:q-setup}
%% ====================================================================

Fix $V = \C^N$ with standard basis $e_1, \ldots, e_N$, as in Part~I.
The tensor power $V^{\otimes n}$ has the same basis $\{e_\mathbf{i}\}$,
but the two group actions of Part~I are replaced by algebraic actions
that deform the $\mathrm{GL}_N$ and $S_n$ symmetries.

\subsection{The quantum group $U_q(\gl_N)$}

\begin{definition}
The \emph{quantum group} $U_q(\gl_N)$ is the $\C$-algebra generated by
$E_i, F_i$ ($1 \leqslant i \leqslant N{-}1$) and $K_j^{\pm 1}$
($1 \leqslant j \leqslant N$) subject to:
\begin{enumerate}[label=(\roman*)]
\item $K_i K_j = K_j K_i$, \; $K_i K_i^{-1} = 1$;
\item $K_i E_j K_i^{-1} = q^{\delta_{ij} - \delta_{i,j+1}} E_j$;
\item $K_i F_j K_i^{-1} = q^{-\delta_{ij} + \delta_{i,j+1}} F_j$;
\item $E_i F_j - F_j E_i = \delta_{ij}
  \dfrac{K_i K_{i+1}^{-1} - K_i^{-1} K_{i+1}}{q - q^{-1}}$;
\item the $q$-Serre relations:
  $E_i^2 E_j - (q + q^{-1}) E_i E_j E_i + E_j E_i^2 = 0$
  for $|i - j| = 1$ (and similarly for $F_i$).
\end{enumerate}
The generators $K_i$ form a \emph{quantum Cartan subalgebra}; $E_i$
and $F_i$ are \emph{quantum raising and lowering operators}.
\end{definition}

At $q = 1$, setting $K_i = 1$ for all $i$ and
$E_i \to E_{i, i+1}$, $F_i \to E_{i+1, i}$, we recover the
universal enveloping algebra $U(\gl_N)$.

\begin{definition}
$U_q(\gl_N)$ acts on $V = \C^N$ by:
\[
  K_j \cdot e_a = q^{\delta_{ja}}\, e_a, \qquad
  E_i \cdot e_a = \delta_{a, i+1}\, e_i, \qquad
  F_i \cdot e_a = \delta_{a, i}\, e_{i+1}.
\]
This is the \emph{vector representation}.
\end{definition}

The action on $V^{\otimes n}$ uses the \emph{coproduct}
$\Delta : U_q(\gl_N) \to U_q(\gl_N) \otimes U_q(\gl_N)$:
\[
  \Delta(K_i) = K_i \otimes K_i, \qquad
  \Delta(E_i) = E_i \otimes K_i K_{i+1}^{-1} + 1 \otimes E_i,
  \qquad
  \Delta(F_i) = F_i \otimes 1 + K_i^{-1} K_{i+1} \otimes F_i.
\]
Iterating $\Delta$ gives an action of $U_q(\gl_N)$ on $V^{\otimes n}$.

\begin{remark}
Unlike the classical case, this action is \emph{not} diagonal: $E_i$
acting on the $k$-th tensor factor inserts a factor of
$K_i K_{i+1}^{-1}$ on all factors to its left. This asymmetry is the
price of deformation, and it is precisely what the $R$-matrix corrects.
\end{remark}

\subsection{The Hecke algebra $\mathrm{H}_n(q)$}

\begin{definition}
The \emph{Iwahori--Hecke algebra} $\mathrm{H}_n(q)$ is the $\C$-algebra
generated by $T_1, \ldots, T_{n-1}$ subject to:
\begin{enumerate}[label=(\roman*)]
\item \emph{Quadratic relation}: $(T_k - q)(T_k + q^{-1}) = 0$,
  equivalently $T_k^2 = (q - q^{-1})\, T_k + 1$;
\item \emph{Braid relations}: $T_k T_{k+1} T_k = T_{k+1} T_k T_{k+1}$;
\item \emph{Far commutativity}: $T_k T_l = T_l T_k$ for $|k - l| > 1$.
\end{enumerate}
\end{definition}

At $q = 1$, the quadratic relation becomes $T_k^2 = 1$, recovering
$\C[S_n]$ with $T_k \mapsto s_k$.

\begin{remark}
The eigenvalues of $T_k$ are $q$ and $-q^{-1}$ (from the quadratic
relation). The dimension of $\mathrm{H}_n(q)$ is $n!$, the same as
$\C[S_n]$, and $\mathrm{H}_n(q) \cong \C[S_n]$ as vector spaces (but
not as algebras for $q \neq 1$).
\end{remark}

\subsection{The $R$-matrix action}

The Hecke algebra acts on $V^{\otimes n}$ via the \emph{$R$-matrix},
which replaces the simple transposition of tensor factors.

\begin{definition}
The \emph{$R$-matrix action} of $T_k$ on $V^{\otimes n}$ acts on
positions $k$ and $k{+}1$ by:
\[
  T_k \cdot (\cdots \otimes e_a \otimes e_b \otimes \cdots)
  = \begin{cases}
    q\, (\cdots \otimes e_a \otimes e_b \otimes \cdots)
      & \text{if } a = b, \\
    (\cdots \otimes e_b \otimes e_a \otimes \cdots)
      & \text{if } a < b, \\
    (\cdots \otimes e_b \otimes e_a \otimes \cdots)
    + (q - q^{-1})\, (\cdots \otimes e_a \otimes e_b \otimes \cdots)
      & \text{if } a > b,
  \end{cases}
\]
where only positions $k$ and $k{+}1$ are affected.
\end{definition}

At $q = 1$ this reduces to the transposition $s_k$ (swapping positions
$k$ and $k{+}1$). For $q \neq 1$, $T_k$ still ``mostly swaps'' but
adds a correction $(q - q^{-1})$ when the larger index is on the left.

\subsection{Commutativity}

\begin{theorem}[Jimbo~\cite{Jimbo}]\label{thm:q-commute}
The actions of $U_q(\gl_N)$ and $\mathrm{H}_n(q)$ on $V^{\otimes n}$
commute: for all $x \in U_q(\gl_N)$ and $h \in \mathrm{H}_n(q)$,
\[
  x \cdot (h \cdot v) = h \cdot (x \cdot v)
  \qquad \text{for all } v \in V^{\otimes n}.
\]
\end{theorem}

\begin{remark}
Unlike the classical case (Proposition~\ref{prop:commute}), this is
\emph{not} obvious. The classical proof relied on $\mathrm{GL}_N$
acting within each factor and $S_n$ permuting factors---these two
operations are manifestly independent. In the quantum case, the
coproduct makes $U_q(\gl_N)$ act asymmetrically across factors, and
the $R$-matrix makes $\mathrm{H}_n(q)$ act non-trivially on values
within each factor. That they still commute is a consequence of the
\emph{Yang--Baxter equation} satisfied by the $R$-matrix.
\end{remark}

%% ====================================================================
\section{The decomposition ($q$-deformed)}
\label{sec:q-decomposition}
%% ====================================================================

\subsection{Partitions and tableaux (recap)}

The combinatorial objects---partitions, Young diagrams, standard Young
tableaux (SYT), and semistandard Young tableaux (SSYT)---are identical
to the classical case (Section~\ref{sec:decomposition}). We use the
same notation: $\lambda \vdash n$ for partitions, $\SYT(\lambda)$ for
standard tableaux, $\SSYT(\lambda, N)$ for semistandard tableaux.

\subsection{The $q$-Schur--Weyl decomposition}

\begin{theorem}[$q$-Schur--Weyl duality]\label{thm:q-SW}
As a $(U_q(\gl_N),\, \mathrm{H}_n(q))$-bimodule,
\begin{equation}\label{eq:q-SW}
  V^{\otimes n}
  \;\cong\;
  \bigoplus_{\substack{\lambda \vdash n \\ \ell(\lambda) \leqslant N}}
  W_q^\lambda \otimes S_q^\lambda,
\end{equation}
where:
\begin{itemize}
\item $W_q^\lambda$ is the irreducible $U_q(\gl_N)$-module of highest
  weight $\lambda$,
\item $S_q^\lambda$ is the irreducible $\mathrm{H}_n(q)$-module indexed
  by $\lambda$.
\end{itemize}
Moreover:
\begin{enumerate}[label=(\alph*)]
\item The natural map
  $\mathrm{H}_n(q) \to \End_{U_q(\gl_N)}(V^{\otimes n})$ is surjective
  (and injective when $n \leqslant N$).
\item The natural map
  $U_q(\gl_N) \to \End_{\mathrm{H}_n(q)}(V^{\otimes n})$ is surjective.
\end{enumerate}
\end{theorem}

\begin{remark}
The dimensions are unchanged:
$\dim W_q^\lambda = \dim W^\lambda = |\SSYT(\lambda, N)|$ and
$\dim S_q^\lambda = \dim S^\lambda = |\SYT(\lambda)|$. The deformation
changes the algebra structure but not the combinatorics.
\end{remark}

%% ====================================================================
\section{Explicit matrix elements for $\mathrm{H}_n(q)$:
  $q$-Young's seminormal form}
\label{sec:q-Young}
%% ====================================================================

\subsection{Content and axial distance (recap)}

The \emph{content} of a box in row $r$, column $c$ of a Young diagram
is $\cont(r,c) = c - r$ (Definition~\ref{def:content}). For a standard
Young tableau $T$, $\cont_T(k)$ is the content of the box containing
$k$. The \emph{axial distance} is $\rho_T(k) = \cont_T(k{+}1) -
\cont_T(k)$ (unchanged from the classical case).

\subsection{The $q$-seminormal representation}

\begin{definition}\label{def:q-Young}
For each $\lambda \vdash n$, the \emph{$q$-seminormal representation}
is the representation of $\mathrm{H}_n(q)$ on the vector space
$\widetilde{S}_q^\lambda = \bigoplus_{T \in \SYT(\lambda)} \C \cdot v_T$
defined by:
\begin{equation}\label{eq:q-Young}
  T_k \cdot v_T = \frac{q - q^{-1}}{1 - q^{-2\rho}}\, v_T
  + \frac{1 - q^{-2\rho - 2}}{1 - q^{-2\rho}}\, v_{s_k \cdot T},
\end{equation}
where $\rho = \rho_T(k)$ and $s_k \cdot T$ is the tableau with
entries $k$ and $k{+}1$ swapped.

If $s_k \cdot T$ is not a valid standard tableau ($|\rho| = 1$), the
second term is absent and
$T_k \cdot v_T = q\, v_T$ (when $\rho = 1$) or
$T_k \cdot v_T = -q^{-1}\, v_T$ (when $\rho = -1$).
\end{definition}

\begin{remark}[Degeneration to $q = 1$]
As $q \to 1$:
\begin{itemize}
\item the diagonal coefficient
  $\dfrac{q - q^{-1}}{1 - q^{-2\rho}} \;\to\; \dfrac{1}{\rho}$
  (by L'H\^opital),
\item the off-diagonal coefficient
  $\dfrac{1 - q^{-2\rho-2}}{1 - q^{-2\rho}} \;\to\;
  \dfrac{\rho + 1}{\rho} = 1 + \dfrac{1}{\rho}$,
\end{itemize}
recovering the classical seminormal form (Remark~3.5).
\end{remark}

\begin{remark}
The representing matrix for $T_k$ in the $2 \times 2$ block
$\{v_T, v_{s_k \cdot T}\}$ is \emph{not} symmetric: the
$(2,1)$-entry uses $\rho' = -\rho$, giving a different off-diagonal
coefficient. This is analogous to the classical seminormal form and
contrasts with the orthogonal form of Part~I where the matrix was
symmetric.
\end{remark}

\begin{example}\label{ex:q-S3-21}
For $\lambda = (2,1) \vdash 3$ with standard tableaux
$T_1 = \;\ytableaushort{12,3}$ and $T_2 = \;\ytableaushort{13,2}$:

The axial distances are $\rho_{T_1}(1) = 1$, $\rho_{T_1}(2) = -2$,
$\rho_{T_2}(1) = -1$, $\rho_{T_2}(2) = 2$.

Since $|\rho_{T_1}(1)| = |\rho_{T_2}(1)| = 1$:
\[
  T_1 \cdot v_{T_1} = q\, v_{T_1}, \qquad
  T_1 \cdot v_{T_2} = -q^{-1}\, v_{T_2}.
\]
For $T_2$ (with $\rho = \rho_{T_1}(2) = -2$):
\[
  T_2 = \begin{pmatrix}
    \dfrac{q - q^{-1}}{1 - q^4} &
    \dfrac{1 - q^{4-2}}{1 - q^4} \\[10pt]
    \dfrac{1 - q^{-4+2}}{1 - q^{-4}} &
    \dfrac{q - q^{-1}}{1 - q^{-4}}
  \end{pmatrix}
  = \begin{pmatrix}
    -\dfrac{q - q^{-1}}{q^4 - 1} &
    \dfrac{1 - q^2}{1 - q^4} \\[10pt]
    \dfrac{1 - q^{-2}}{1 - q^{-4}} &
    \dfrac{q - q^{-1}}{1 - q^{-4}}
  \end{pmatrix}.
\]
One can verify $(T_1 - q)(T_1 + q^{-1}) = 0$,
$(T_2 - q)(T_2 + q^{-1}) = 0$, and $T_1 T_2 T_1 = T_2 T_1 T_2$.
\end{example}

%% ====================================================================
\section{Explicit matrix elements for $U_q(\gl_N)$: the $q$-GT basis}
\label{sec:q-GT}
%% ====================================================================

\subsection{$q$-GT patterns (recap)}

Gel'fand--Tsetlin patterns are unchanged: a GT pattern $\Lambda$ of
shape $\lambda$ is a triangular array satisfying the interlacing
conditions $\lambda_{k,i} \geqslant \lambda_{k-1,i} \geqslant
\lambda_{k,i+1}$ (Definition~4.1). The bijection between GT patterns
and semistandard tableaux is the same.

\subsection{The $U_q(\gl_N)$ action in the $q$-GT basis}

Write $v_\Lambda$ for the basis vector of $W_q^\lambda$ indexed by
$\Lambda$. In the shifted coordinates $\ell_{k,i} = \lambda_{k,i} - i$
(Remark~\ref{rem:shifted}), the action of $U_q(\gl_N)$ is:

\begin{proposition}[$q$-GT formulae]\label{prop:q-GT}
\begin{align}
  K_k\, v_\Lambda
  &= q^{\sum_{i=1}^k \lambda_{k,i} - \sum_{i=1}^{k-1}
     \lambda_{k-1,i}}\, v_\Lambda,
  \label{eq:q-GT-K} \\[6pt]
  E_k\, v_\Lambda
  &= -\sum_{i=1}^k
    \frac{\displaystyle\prod_{j=1}^{k+1}[\ell_{k,i} - \ell_{k+1,j}]_q
    \;\prod_{j=1}^{k-1}[\ell_{k,i} - \ell_{k-1,j}]_q}
    {\displaystyle\prod_{\substack{j=1 \\ j \neq i}}^{k}
    [\ell_{k,i} - \ell_{k,j}]_q}
    \; v_{\Lambda + \delta_{k,i}},
  \label{eq:q-GT-E} \\[6pt]
  F_k\, v_\Lambda
  &= \sum_{i=1}^k
    \frac{\displaystyle\prod_{j=1}^{k+1}[\ell_{k,i} - \ell_{k+1,j}]_q
    \;\prod_{j=1}^{k-1}[\ell_{k,i} - \ell_{k-1,j}]_q}
    {\displaystyle\prod_{\substack{j=1 \\ j \neq i}}^{k}
    [\ell_{k,i} - \ell_{k,j}]_q}
    \; v_{\Lambda - \delta_{k,i}},
  \label{eq:q-GT-F}
\end{align}
where $[x]_q = (q^x - q^{-x})/(q - q^{-1})$ and the term is zero if
the result is not a valid GT pattern.
\end{proposition}

\begin{remark}
Comparing with the classical GT formulae
(\eqref{eq:GT-raise-shifted}--\eqref{eq:GT-lower-shifted}), the only
change is that each linear factor $(\ell_{k,i} - \ell_{k',j})$ is
replaced by its $q$-analogue $[\ell_{k,i} - \ell_{k',j}]_q$. The
structure---products over adjacent rows in the numerator, Vandermonde
product within the same row in the denominator---is identical.

The Cartan action~\eqref{eq:q-GT-K} exponentiates the classical
weight: $K_k$ has eigenvalue $q^{w_k}$ where $w_k$ is the classical
$E_{k,k}$ eigenvalue~\eqref{eq:GT-diag}.
\end{remark}

%% ====================================================================
\section{The RSK correspondence ($q$-deformed setting)}
\label{sec:q-RSK}
%% ====================================================================

The RSK correspondence is purely combinatorial and does not involve
the parameter $q$. The bijection
\[
  \mathrm{RSK} : \{1, \ldots, N\}^n
  \;\xrightarrow{\;\sim\;}\;
  \bigsqcup_{\substack{\lambda \vdash n \\ \ell(\lambda) \leqslant N}}
  \SSYT(\lambda, N) \times \SYT(\lambda)
\]
is identical to Definition~5.1. The insertion algorithm, recording
tableau, and all combinatorial properties (Proposition~\ref{prop:RSK-props})
are unchanged.

%% ====================================================================
\section{RSK is the $q$-Schur--Weyl decomposition}
\label{sec:q-RSK-is-SW}
%% ====================================================================

The interpretive result of Section~\ref{sec:RSK-is-SW} carries over
verbatim to the $q$-deformed setting.

\begin{theorem}[$q$-RSK as Schur--Weyl]\label{thm:q-RSK-SW}
Under the isomorphism~\eqref{eq:q-SW}, the basis element
$e_\mathbf{i}$ corresponds to
\[
  e_\mathbf{i} \;\longleftrightarrow\;
  w_P \otimes v_Q
  \;\in\; W_q^\lambda \otimes S_q^\lambda,
\]
where $(P, Q) = \mathrm{RSK}(\mathbf{i})$. The insertion tableau $P$
indexes the $U_q(\gl_N)$-irrep basis (via the $q$-GT identification),
and the recording tableau $Q$ indexes the $\mathrm{H}_n(q)$-irrep
basis (via the SYT basis of the $q$-seminormal representation).
\end{theorem}

\begin{remark}
The table of Section~6.2 applies without change: $P$ is sensitive to
what letters appear ($U_q(\gl_N)$ symmetry), $Q$ is sensitive to
the order of arrival ($\mathrm{H}_n(q)$ symmetry), and the shape
$\lambda$ determines which summand.
\end{remark}

%% ====================================================================
\section{The bridge: $q$-Jucys--Murphy elements}
\label{sec:q-bridge}
%% ====================================================================

\subsection{Motivation}

The strategy of Part~I carries over: use the chain
$\mathrm{H}_1(q) \subset \mathrm{H}_2(q) \subset \cdots \subset
\mathrm{H}_n(q)$ to build a maximal commutative subalgebra whose
spectral projectors carve out the irreducible summands.

In the classical case, $J_k = X_k - X_{k-1}$ was an additive
``peeling off'' of the new element. In the $q$-deformed case, the
construction becomes multiplicative.

\subsection{The $q$-Jucys--Murphy elements}

\begin{definition}\label{def:q-JM}
The \emph{$q$-Jucys--Murphy elements} of $\mathrm{H}_n(q)$ are
\[
  L_1 = 1, \qquad
  L_k = T_{k-1}\, L_{k-1}\, T_{k-1},
  \quad k = 2, \ldots, n.
\]
\end{definition}

Expanding the recursion: $L_2 = T_1^2 = (q - q^{-1}) T_1 + 1$,
$L_3 = T_2 L_2 T_2 = T_2 T_1^2 T_2$, and so on. Each $L_k$ is a
product of Hecke generators involving $T_1, \ldots, T_{k-1}$.

\begin{remark}[Degeneration to $q = 1$]
At $q = 1$, the recursion gives $L_k = s_{k-1} L_{k-1} s_{k-1}$. From
the classical recursion $J_k = s_{k-1} J_{k-1} s_{k-1} + s_{k-1}$
(Lemma~\ref{lem:JM-recursion}), we see that $L_k$ and $J_k$ are
related but not identical at $q = 1$: the $q$-JM elements are
\emph{multiplicative} while the classical ones are \emph{additive}.
The precise relation is $L_k|_{q=1} = J_k + 1$ (the identity shift
absorbs the $+s_{k-1}$ correction).
\end{remark}

\subsection{Commutativity}

\begin{lemma}\label{lem:q-JM-commute}
The elements $L_1, \ldots, L_n$ commute pairwise.
\end{lemma}

\begin{proof}
We first show that $L_k$ commutes with $T_i$ for $i \neq k{-}1, k$.

For $i < k{-}1$: by induction on $k$. The base case $L_2 =
(q-q^{-1})T_1 + 1$ trivially commutes with $T_i$ for $i > 2$ (far
commutativity). For the inductive step,
$L_k = T_{k-1} L_{k-1} T_{k-1}$, and $T_i$ commutes with $T_{k-1}$
(since $i < k{-}1$ implies $|i - (k{-}1)| > 1$) and with $L_{k-1}$
(by induction), hence with $L_k$.

For $i > k$: $T_i$ commutes with all of $T_1, \ldots, T_{k-1}$, hence
with $L_k$.

Now $L_k$ and $L_l$ commute for $|k - l| \geqslant 2$:
$L_l = T_{l-1} L_{l-1} T_{l-1}$ involves only $T_{l-1}$ and
$L_{l-1}$, both of which commute with $L_k$ (the former because
$l - 1 > k$, the latter by induction).

For $L_k$ and $L_{k+1}$: $L_{k+1} = T_k L_k T_k$, so
$L_k L_{k+1} = L_k T_k L_k T_k$ and
$L_{k+1} L_k = T_k L_k T_k L_k$. These are equal because
$L_k T_k L_k = T_k L_k T_k \cdot L_k \cdot T_k^{-1}$...
more precisely, one verifies directly from the Hecke relations and
the recursion that
$T_k L_k T_k \cdot L_k = L_k \cdot T_k L_k T_k$.
\end{proof}

The $L_k$ generate the \emph{$q$-Gelfand--Tsetlin subalgebra} of
$\mathrm{H}_n(q)$. By ``commutative subalgebra'' we mean that the
$L_k$ commute with each other---they do not commute with arbitrary
elements of $\mathrm{H}_n(q)$.

\subsection{The eigenvalue theorem}

\begin{proposition}\label{prop:q-JM-eigenvalues}
In the $q$-seminormal representation $\widetilde{S}_q^\lambda$
(Definition~\ref{def:q-Young}), the $q$-Jucys--Murphy element $L_k$
acts diagonally in the SYT basis:
\[
  L_k \cdot v_T = q^{2\,\cont_T(k)}\, v_T,
\]
where $\cont_T(k) = c - r$ is the content of the box in row~$r$,
column~$c$ that contains entry~$k$ in~$T$.
\end{proposition}

\begin{proof}
By induction on $k$, using the recursion
$L_k = T_{k-1} L_{k-1} T_{k-1}$.

\emph{Base case ($k = 1$):} $L_1 = 1$, and $q^{2 \cdot 0} = 1$
(entry $1$ is always in box $(1,1)$ with content $0$).

\emph{Inductive step:} Assume $L_{k-1}\, v_T = q^{2 c_1}\, v_T$
for all $T$, where $c_1 = \cont_T(k{-}1)$. Set $c_2 = \cont_T(k)$,
$\rho = c_2 - c_1$, and let $T' = s_{k-1} \cdot T$. We must show
$L_k\, v_T = T_{k-1} L_{k-1} T_{k-1}\, v_T = q^{2 c_2}\, v_T$.

Write $\alpha = (q - q^{-1})/(1 - q^{-2\rho})$ and
$\beta = (1 - q^{-2\rho-2})/(1 - q^{-2\rho})$ for the diagonal and
off-diagonal coefficients in~\eqref{eq:q-Young}, and $\alpha' =
(q - q^{-1})/(1 - q^{2\rho})$, $\beta' = (1 - q^{2\rho-2})/(1 -
q^{2\rho})$ for the corresponding coefficients for $T'$ (where
$\rho_{T'}(k{-}1) = -\rho$). Compute step by step:
\begin{align*}
  T_{k-1}\, v_T &= \alpha\, v_T + \beta\, v_{T'}, \\
  L_{k-1}(\alpha\, v_T + \beta\, v_{T'})
    &= \alpha\, q^{2c_1}\, v_T + \beta\, q^{2c_2}\, v_{T'}
    \quad\text{(by induction, since $\cont_{T'}(k{-}1) = c_2$)}, \\
  T_{k-1}(\alpha\, q^{2c_1}\, v_T + \beta\, q^{2c_2}\, v_{T'})
    &= \alpha q^{2c_1}(\alpha\, v_T + \beta\, v_{T'})
    + \beta q^{2c_2}(\beta'\, v_T + \alpha'\, v_{T'}).
\end{align*}
The $v_T$ coefficient is
$\alpha^2 q^{2c_1} + \beta \beta' q^{2c_2}$.
One verifies (using $q^{2c_2} = q^{2\rho} q^{2c_1}$ and the
explicit formulae for $\alpha, \beta, \beta'$) that this equals
$q^{2c_2}$.

The $v_{T'}$ coefficient is
$\alpha \beta q^{2c_1} + \beta \alpha' q^{2c_2}
= \beta(\alpha q^{2c_1} + \alpha' q^{2c_2})$.
Since $\alpha + \alpha' = q - q^{-1}$ and
$\alpha q^{2c_1} + \alpha' q^{2c_2}
= \alpha q^{2c_1} + (q - q^{-1} - \alpha) q^{2\rho} q^{2c_1}
= q^{2c_1}[\alpha(1 - q^{2\rho}) + (q-q^{-1})q^{2\rho}]$,
and substituting $\alpha(1 - q^{2\rho}) = -(q-q^{-1})q^{2\rho}$
(from the definition of $\alpha$), this is $0$.

So $L_k\, v_T = q^{2c_2}\, v_T = q^{2\,\cont_T(k)}\, v_T$.
\end{proof}

\begin{remark}
At $q = 1$: $L_k|_{q=1} = 1$ for all $k$ (since $L_k = T_{k-1}
L_{k-1} T_{k-1}$ and $T_k|_{q=1} = s_k$ with $s_k^2 = 1$, giving
$L_k = s_{k-1} L_{k-1} s_{k-1}$ with $L_1 = 1$, so $L_k = 1$ for
all $k$). This is consistent with $q^{2\,\cont_T(k)}|_{q=1} = 1$.
The classical JM elements $J_k$ (with non-trivial eigenvalues
$\cont_T(k)$) are recovered via the infinitesimal version:
$J_k = \lim_{q \to 1} \frac{L_k - 1}{q^2 - 1}$.
\end{remark}

\subsection{Content separation and maximality}

Lemma~\ref{lem:content-separates} (different SYTs have different
content vectors) applies unchanged, since content vectors do not
depend on $q$. The $q$-eigenvalues $q^{2\,\cont_T(k)}$ are distinct
whenever the contents are distinct (since $q$ is not a root of unity),
so the $q$-GZ subalgebra separates all basis vectors.

\begin{theorem}\label{thm:q-GZ-maximal}
The $q$-GZ subalgebra generated by $L_1, \ldots, L_n$ is a maximal
commutative subalgebra of $\mathrm{H}_n(q)$.
\end{theorem}

\begin{proof}
The same argument as Theorem~\ref{thm:GZ-maximal}: by
Artin--Wedderburn, $\mathrm{H}_n(q) \cong
\bigoplus_\lambda \mathrm{Mat}_{d_\lambda}(\C)$ (since $q$ is not a
root of unity), the $q$-GZ subalgebra acts diagonally in the SYT
basis with distinct eigenvalues, hence contains all diagonal matrices
in each block.
\end{proof}

\subsection{$q$-spectral projectors}

\begin{definition}\label{def:q-primitive-idempotent}
For each $Q \in \SYT(\lambda)$, the \emph{$q$-primitive idempotent} is
\[
  e_Q^{(q)} \;=\; \prod_{k=2}^{n}\;\prod_{\substack{c \neq
    q^{2\,\cont_Q(k)} \\
    c \text{ an eigenvalue of } L_k}}
    \frac{L_k - c}{q^{2\,\cont_Q(k)} - c}
  \;\;\in\; \mathrm{H}_n(q).
\]
\end{definition}

Since $e_Q^{(q)} \in \mathrm{H}_n(q)$ and the $U_q(\gl_N)$-action
commutes with $\mathrm{H}_n(q)$
(Theorem~\ref{thm:q-commute}), the subspace
$e_Q^{(q)} \cdot V^{\otimes n}$ is an irreducible
$U_q(\gl_N)$-module isomorphic to $W_q^\lambda$, exactly as in
the classical case (Proposition~\ref{prop:JM-bridge}).

%% ====================================================================
\section{Degeneration: from $q$ to $1$}
\label{sec:degeneration}
%% ====================================================================

\begin{center}
\renewcommand{\arraystretch}{1.5}
\begin{tabular}{c|c|c}
  & \textbf{Classical ($q = 1$)}
  & \textbf{$q$-deformed} \\
\hline
\textbf{Symmetry algebra}
  & $\C[S_n]$ & $\mathrm{H}_n(q)$ \\
\textbf{Generators}
  & $s_k$ with $s_k^2 = 1$
  & $T_k$ with $(T_k - q)(T_k + q^{-1}) = 0$ \\
\textbf{Dual algebra}
  & $U(\gl_N)$ & $U_q(\gl_N)$ \\
\textbf{Action on $V^{\otimes n}$}
  & diagonal + permutation
  & coproduct + $R$-matrix \\
\textbf{JM elements}
  & $J_k = \textstyle\sum_{i<k} (i\;k)$ (additive)
  & $L_k = T_{k-1} L_{k-1} T_{k-1}$ (multiplicative) \\
\textbf{JM recursion}
  & $J_k = s_{k-1} J_{k-1} s_{k-1} + s_{k-1}$
  & $L_k = T_{k-1} L_{k-1} T_{k-1}$ \\
\textbf{JM eigenvalues}
  & $\cont_T(k) \in \Z$
  & $q^{2\,\cont_T(k)} \in \C^*$ \\
\textbf{GT coefficients}
  & $(\ell_{k,i} - \ell_{k',j})$
  & $[\ell_{k,i} - \ell_{k',j}]_q$ \\
\textbf{Cartan eigenvalue}
  & $\sum \lambda_{k,i} - \sum \lambda_{k-1,i}$
  & $q^{\sum \lambda_{k,i} - \sum \lambda_{k-1,i}}$ \\
\textbf{Combinatorics}
  & \multicolumn{2}{c}{unchanged: SYT, SSYT, GT patterns, RSK}
\end{tabular}
\end{center}

The deformation is systematic: \emph{integers become $q$-integers,
contents become $q$-contents, linear factors become $q$-linear
factors}. The combinatorial backbone---partitions, tableaux, RSK---is
untouched. The algebraic structures deform continuously, and the
representation theory deforms with them.

\begin{remark}
The $q$-JM recursion $L_k = T_{k-1} L_{k-1} T_{k-1}$ is
\emph{simpler} than the classical one
$J_k = s_{k-1} J_{k-1} s_{k-1} + s_{k-1}$. The additive correction
$+s_{k-1}$ in the classical case is an artifact of the degeneration
$q \to 1$: the quadratic relation $(T_k - q)(T_k + q^{-1}) = 0$
absorbs this correction, and it only reappears when the quadratic
degenerates to $s_k^2 = 1$.
\end{remark}

\begin{thebibliography}{99}

\bibitem{DM}
A.~Davydov and A.~Molev,
\emph{A categorical approach to classical and quantum Schur--Weyl duality},
preprint.

\bibitem{Fulton}
W.~Fulton,
\emph{Young Tableaux},
London Mathematical Society Student Texts 35,
Cambridge University Press, 1997.

\bibitem{GoodmanWallach}
R.~Goodman and N.~R.~Wallach,
\emph{Symmetry, Representations, and Invariants},
Graduate Texts in Mathematics 255, Springer, 2009.

\bibitem{Jimbo}
M.~Jimbo,
\emph{A $q$-analogue of $U(\gl(N+1))$, Hecke algebra, and the
Yang--Baxter equation},
Lett.\ Math.\ Phys.\ \textbf{11} (1986), 247--252.

\bibitem{Knuth}
D.~E.~Knuth,
\emph{Permutations, matrices, and generalized Young tableaux},
Pacific J.\ Math.\ \textbf{34} (1970), 709--727.

\bibitem{Mathas}
A.~Mathas,
\emph{Iwahori--Hecke Algebras and Schur Algebras of the Symmetric
Group},
University Lecture Series 15,
American Mathematical Society, 1999.

\bibitem{Molev}
A.~Molev,
\emph{Yangians and Classical Lie Algebras},
Mathematical Surveys and Monographs 143,
American Mathematical Society, 2007.

\bibitem{OkounkovVershik}
A.~Okounkov and A.~Vershik,
\emph{A new approach to representation theory of symmetric groups},
Selecta Math.\ (N.S.) \textbf{2} (1996), 581--605.

\bibitem{Ram}
A.~Ram,
\emph{Seminormal representations of Weyl groups and Iwahori--Hecke
algebras},
Proc.\ London Math.\ Soc.\ (3) \textbf{75} (1997), 99--133.

\bibitem{Sagan}
B.~E.~Sagan,
\emph{The Symmetric Group: Representations, Combinatorial Algorithms,
and Symmetric Functions}, 2nd ed.,
Graduate Texts in Mathematics 203, Springer, 2001.

\bibitem{Stanley}
R.~P.~Stanley,
\emph{Enumerative Combinatorics}, vol.~2,
Cambridge Studies in Advanced Mathematics 62,
Cambridge University Press, 1999.

\end{thebibliography}

\end{document}
